% Auto-generated from meeting_2569-05-09_admission2569-interview.md (กบ.วช. format v2)
\documentclass{meetingminutes}
\meetingcommittee{คณะกรรมการบริหารหลักสูตร วิศวกรรมคอมพิวเตอร์และปัญญาประดิษฐ์ (บัณฑิตศึกษา)}
\meetingnumber{๑}{๒๕๖๙}
\meetingdate{วันที่ ๙ พฤษภาคม พ.ศ. ๒๕๖๙}
\meetingplace{ห้องประชุมคณะเทคโนโลยีอุตสาหกรรม}
\meetingstart{๐๙.๐๐ น.}
\meetingend{๑๔.๓๐ น.}

\begin{document}
\maketitle

\psection{ผู้มาประชุม}
\begin{participants}
\end{participants}

\psection{ผู้ไม่มาประชุม}
\noindent ไม่มี\par

\startmeeting
\openingchair{ผู้ช่วยศาสตราจารย์ ดร.พัชรี มณีรัตน์}{กล่าวเปิดการประชุมและดำเนินการประชุมตามระเบียบวาระ ดังนี้}

\agenda{๑}{รายงานจำนวนผู้สมัคร ปีการศึกษา 2569 ภาคที่ 1}
ฝ่ายเลขานุการรายงานว่ามีผู้สมัครเข้าศึกษา ปีการศึกษา 2569 ภาคที่ 1 จำนวนทั้งสิ้น 8 คน ซึ่งเพิ่มขึ้นจากปีการศึกษา 2568 อย่างมีนัยสำคัญ (133\% เพิ่มขึ้น) แสดงถึงการรับรู้หลักสูตรในกลุ่มเป้าหมายที่ดีขึ้น ผู้สมัครทั้ง 8 คนผ่านการตรวจสอบคุณสมบัติเบื้องต้นและได้รับการนัดหมายสอบสัมภาษณ์


\agenda{๒}{ดำเนินการสอบสัมภาษณ์ผู้สมัคร 8 คน}
คณะกรรมการดำเนินการสอบสัมภาษณ์ผู้สมัครทั้ง 8 คน โดยใช้เกณฑ์เดียวกับปีการศึกษา 2568 ประกอบด้วย: 1. ความรู้พื้นฐานด้านวิศวกรรมคอมพิวเตอร์/AI 2. แรงจูงใจและเป้าหมายการศึกษา 3. แผนการวิจัย/สารนิพนธ์เบื้องต้น 4. ความสามารถสื่อสารและแสดงความคิดเห็น


\agenda{๓}{สรุปผลการสัมภาษณ์}
คณะกรรมการประเมินและสรุปผลการสัมภาษณ์ผู้สมัครทั้ง 8 คน ผลปรากฏว่าผู้สมัครทุกคนมีคุณสมบัติและศักยภาพที่เหมาะสม

\resolution{1. รับผู้สมัครทั้ง 8 คนเข้าศึกษา โดยรอยืนยันการรายงานตัวอีกครั้ง
2. กำหนดรายงานตัวและขึ้นทะเบียนนักศึกษา: ต้นเดือนมิถุนายน 2569
3. เปิดภาค 1/2569: 24 มิถุนายน 2569}

\agenda{๔}{เรื่องอื่น ๆ}
ที่ประชุมรับทราบว่า ผลสำรวจความพึงพอใจนักศึกษาปัจจุบัน (n=4) ได้ดำเนินการแล้ว คะแนนเฉลี่ยรวม 4.94/5.00 ซึ่งสะท้อนความพึงพอใจในระดับสูงมาก และจะนำเสนอในการประชุมครั้งต่อไป


\adjournmeeting

\begin{sigblock}
\typist{ผู้ช่วยศาสตราจารย์ ดร.กาญจนา ดาวเด่น}
\reviewer{ผู้ช่วยศาสตราจารย์ ดร.พัชรี มณีรัตน์}{กรรมการ}
\chairsig{ผู้ช่วยศาสตราจารย์ ดร.พัชรี มณีรัตน์}{ประธานการประชุม}
\end{sigblock}

\end{document}
